\subsection{3.3 Aufzeichnungen müssen die richtige Interferenz
erhalten}\label{aufzeichnungen-muxfcssen-die-richtige-interferenz-erhalten}

Für den antisymmetrischen Paaradapter

\[K:\mathbb C^4\otimes\mathbb C^4\longrightarrow\Lambda^2\mathbb C^4\]

gilt in der hier verwendeten Normierung

\[K^\dagger K=I-\mathsf S,\]

mit dem Swap \(\mathsf S\). Die Wege \(ab\) und \(ba\) treffen denselben
Vermittler mit entgegengesetztem Vorzeichen. Fügt man orthogonale
Records an diese zwei Wege an, verschwinden ihre Kreuzterme. Man erhält
stattdessen \(I-D_{\mathrm{diag}}\), wobei
\(D_{\mathrm{diag}}=\sum_a|aa\rangle\langle aa|\).

Damit verändert eine unpassende Aufzeichnung die Dynamik. Kanten dürfen
als verschiedene Ereignisse unterscheidbar bleiben; vertauschte innere
Wege derselben antisymmetrischen Paarung müssen kohärent bleiben. Die
Aussage ist nicht „History ist unmöglich``, sondern eine konkrete
Verträglichkeitsbedingung an ihre Gramstruktur. {[}S072, Nachtrag N3{]}

Allgemein gilt für zwei lineare Adapter desselben Eingangsraums:

\[L^\dagger L=K^\dagger K
\quad\Longleftrightarrow\quad L=VK\]

mit einer Isometrie \(V\) auf \(\operatorname{ran}K\). Der Beweis steht
in Abschnitt 8. Diese einfache Identität wird zum ersten Bestandteil
einer gemeinsamen Rekonstruktion.

\section{4. Die Vierträgerzelle und die Grenze lokaler
Modelllösungen}\label{die-viertruxe4gerzelle-und-die-grenze-lokaler-modellluxf6sungen}

Der total antisymmetrische Zustand

\[|\Omega_4\rangle=\frac1{\sqrt{24}}
\sum_{\pi\in S_4}\operatorname{sgn}(\pi)
|\pi(0)\pi(1)\pi(2)\pi(3)\rangle\]

ist ein \(SU(4)\)-Singulett. Seine Einteilchenreduktionen sind
\(I_4/4\), seine Paarreduktionen \((I-\mathsf S_{ij})/12\). Die Ordnung
liegt in den Beziehungen, obwohl jedes einzelne Teil maximal gemischt
aussieht.

Für

\[H_{\mathrm{tet}}=J\sum_{i<j}P^+_{ij},\qquad
P^+_{ij}=\tfrac12(I+\mathsf S_{ij}),\quad J>0\]

ist \(\Omega_4\) der eindeutige Nullzustand. Fordert man dieselbe
vollständige Antisymmetrie bei mehr als vier Viererträgern, ist
\(\Lambda^n\mathbb C^4=0\) für \(n>4\). Eine globale Welt kann deshalb
nicht aus beliebig vielen gleichzeitig perfekt erfüllten solchen
Paarbedingungen entstehen. Frustration, andere Verklebung oder andere
Freiheitsgrade werden wesentlich. {[}S046, Kap. 6--8{]}

Für zwei vollständige Tetramer mit der festgelegten positiven
Einzelbrücke existiert ein allkoppliger Grundzustandssatz. Mit
\(\lambda\ge0\) und

\[R=\sqrt{16J^2-2J\lambda+\lambda^2},\qquad
Q=\sqrt{4J^2+\lambda^2}\]

gilt

\[E_0=\tfrac12(4J+\lambda-R),\qquad
\operatorname{gap}=J+\tfrac12(R-Q)>J/2.\]

Historische Schwellen \(\lambda=8J\) waren Grenzen schwächerer Beweise,
keine Phasenübergänge dieses Modells. Dieser Satz betrifft zwei Zellen
und die bezeichnete Brücke. Er ist weder eine uniforme Vielzellenaussage
noch der Grundzustandssatz der nativen 64/60-Bank.

Ein Produkt eindimensionaler Singulettgrundräume ist selbst
eindimensional. Projiziert man vollständig darauf, bleibt nur eine
skalare Energie. Nichttriviale niedrige Dynamik benötigt Anregungen,
Randfreiheiten oder Multiplizitäten. Diese Beobachtung verbindet das
Zellmodell mit der späteren Frage, wo Raum überhaupt mathematisch Platz
haben kann.

\section{5. Die native Bank: festgelegte Dynamik und gesicherte
Antwort}\label{die-native-bank-festgelegte-dynamik-und-gesicherte-antwort}

\subsection{5.1 Modellvertrag und
Tensor}\label{modellvertrag-und-tensor}

Die native Konstruktion verwendet Fermionoperatoren \(f_i\) für
\(i=1,\dots,64\), Bosonoperatoren \(b_A\) für \(A=1,\dots,60\) und

\[
P_A=\sum_{i<j}W_{A,ij}f_jf_i,\qquad
H_{\mathrm{nat}}=\Delta N_b+g\sum_A(b_A^\dagger P_A+P_A^\dagger b_A).
\]

Es gelten \(\Delta>0\), reelles \(g\) und

\[N=N_f+2N_b,\qquad [H_{\mathrm{nat}},N]=0.\]

\(W\) hat 480 Einträge \(\pm1\), acht disjunkte Paare je Zeile und
\(WW^\dagger=8I_{60}\). Die Fermionen tragen \((16,4)\), die Bosonen
\((10,6)\) unter der inneren Gruppe. Der verwendete Tensorpin lautet:

\texttt{3f00a0892157a691e4ef26920c3702772c7bf76a250fad98ff04a079bb64b763}.

Der Tensor und seine Gramidentität wurden in dieser Arbeit frisch
geprüft. Die gesamte Fockbank ist wegen der Bosonen
unendlichdimensional; jeder feste nichtnegative \(N\)-Sektor ist
endlichdimensional. „Endliche Bank`` bedeutet daher eine endliche Zahl
von Moden, nicht automatisch einen endlichen Gesamthilbertraum.

\subsection{5.2 Grundzustand und
Entnahmelinie}\label{grundzustand-und-entnahmelinie}

Der übernommene native Satz gilt für den \(\mu=0\)-Vertrag und
\(0<|g|/\Delta\le1/20\): Der globale Grundzustand ist eindeutig und
liegt bei \(N=64\). Am Betriebspunkt \(g/\Delta=1/20\) geben die
bisherigen Zertifikate

\[-1.158089<E_0/\Delta<-1.129636,
\qquad 0.842846<\langle N_b\rangle<1.245656.\]

Die neue v1.6.7-Kompression verschärft die Obergrenze zu

\[\boxed{-1.158089<E_0/\Delta<-1.12963811.}\]

Für die isolierte niedrige Entnahmelinie im \(N=63\)-Sektor gilt damit

\[\boxed{0.00773911<\epsilon/\Delta<0.039079764,}\]

sowie das übernommene Gewicht

\[Z_{\mathrm{low}}>\frac{40912436089}{46487375000}>0.88007628.\]

Die Normierung ist wesentlich: \(Z_{\mathrm{low}}\) bezieht sich auf das
gesamte normierte Fermionspektralgewicht je Mode. Die gesamte
Entnahmenorm ist dagegen
\(\nu=\langle f_r^\dagger f_r\rangle=1-\langle N_b\rangle/32\). Ein
isolierter Pol im endlichen Sektor ist zunächst eine isolierte
Eigenlinie; seine Identifikation als relativistisches Teilchen verlangt
Dispersion, Feldadapter und Grenzwert. {[}S033, §1; S068, §4{]}

\subsection{5.3 Weshalb eine schwächere Sonde den stärkeren Satz nicht
widerlegt}\label{weshalb-eine-schwuxe4chere-sonde-den-stuxe4rkeren-satz-nicht-widerlegt}

Ein paralleler Arbeitsstand konnte den Grundzustand mit seinen eigenen
gröberen Schranken bei \(g/\Delta=1/20\) nicht isolieren. Er verlangte
dazu Energieobergrenzen, die mit dem stärkeren früheren Intervall
unvereinbar wären. Daraus folgt, dass diese Sonde zu schwach ist. Es
folgt weder eine neue niedrigere Energie noch die Widerlegung des
vorhandenen nativen Satzes. Ebenso ist ein Ritzwert oberhalb einer
bekannten besseren Obergrenze keine neue Grundenergie.

Versionskonflikte betreffen auch Prüfberichte: In einer früheren Datei
war eine Nichtschlussnorm um \(229^2\) falsch skaliert. Die aktuellen
Korrekturen verwenden

\[\|w_2\|^2=\frac{5001523200}{229}>0.\]

Solche Unterschiede werden durch Dateipins und direkte Rechnung
entschieden, nicht durch das Wort PASS im Bericht. {[}S033, §2; S068,
§2{]}

\subsection{5.4 Zustandswahl ist eine zusätzliche physikalische
Frage}\label{zustandswahl-ist-eine-zusuxe4tzliche-physikalische-frage}

Die Familie

\[H_\mu=H_{\mathrm{nat}}+\mu N\]

bewahrt dieselbe innere Symmetrie. Für neutrale \(O\) mit \([O,N]=0\)
gilt bei demselben Anfangszustand exakt

\[e^{itH_\mu}Oe^{-itH_\mu}=e^{itH_{\mathrm{nat}}}Oe^{-itH_{\mathrm{nat}}}.\]

Dennoch verändert \(\mu\) den Vergleich verschiedener Ladungssektoren.
Die dokumentierte Schranke bei \(\mu=\Delta/50\) und \(g=\Delta/20\)
wählt das leere Vakuum. Symmetrie und neutrale Daten allein bestimmen
somit den globalen Zustandsvertrag nicht. In einem festgehaltenen
Gesamtsektor ist \(\mu N\) dagegen nur eine Konstante und kein
beobachtbarer Unterschied. Eine fundamentale Theorie benötigt Auswahl
nur bis auf Gleichheit aller zugelassenen Prozesse. {[}S046, Kap. 36;
S003, §6{]}

\subsection{5.5 Der native Tensor als verlustfreier
Kodierer}\label{der-native-tensor-als-verlustfreier-kodierer}

Der während der Synthese hinzugekommene Nachtrag zu v1.6.7 macht eine
weitere Verbindung explizit. Aus derselben Gramidentität folgt

\[V=\frac{W^\dagger}{\sqrt8}:\mathbb C^{60}\longrightarrow\Lambda^2\mathbb C^{64},
\qquad V^\dagger V=I_{60},\qquad
\Pi=VV^\dagger=\frac{W^\dagger W}{8}.\]

Ein logischer Zustand mit 60 Komponenten hat damit eine verlustfreie
Darstellung als kohärenter Fermionpaarzustand. Die Rückabbildung sieht
jedoch nur einen Rang-60-Teilraum der insgesamt 2016 Paarrichtungen; ihr
Kern ist 1956-dimensional. Diese dunklen Richtungen sind für diese
Abbildung unsichtbar, nicht allgemein bedeutungslos. Im hellen
\(N=2\)-Bereich lautet die gemeinsame Dynamik exakt

\[H_{2,\mathrm{hell}}=
\begin{pmatrix}0&\sqrt8g\\\sqrt8g&\Delta\end{pmatrix}\otimes I_{60}.\]

Sie wandelt Paar- und Bosondarstellung ineinander um. Bei
\(g/\Delta=1/20\) beträgt die maximale Umwandlungswahrscheinlichkeit
\(2/27\). Der allein in die Paarseite eingebettete Code ist für
\(g\ne0\) kein invarianter Hamiltonunterraum. Dies ist eine konkrete
Verbindung zweier Ansichten desselben \(N=2\)-Prozesses; sie
identifiziert diesen Prozess nicht mit dem Grundzustand in \(N=64\).
{[}S074, B5.1{]}

\textbf{Ein präziser Schutztest.} Für jede einzelne Fermionmode ergibt
sich

\[\operatorname{spec}(V^\dagger n_rV)
=\{0\text{ mit Vielfachheit }45,\quad
\tfrac18\text{ mit Vielfachheit }15\}.\]

Die Besetzung ist auf dem logischen Raum nicht skalar. Eine verlorene
Mode kann daher Information über den kodierten Zustand an die Umgebung
geben. Der gesamte 60-dimensionale Code korrigiert die beliebige
Löschung einer einzelnen Mode nicht. Die Kompressionen aller 64 Moden
wurden für dieses Paper am eingefrorenen Tensor erneut exakt geprüft.
Kleinere Untercodes und kontrolliert zusammengesetzte Kodierungsnetze
bleiben mögliche Forschungsrichtungen. Isometrie allein begründet weder
Fehlerkorrektur noch eine holografische Raumgrenze. Etablierte
holografische Codes liefern dafür zusätzliche Rekonstruktions- und
Schutzbedingungen. {[}S074, B5.2; E11; E12{]}

\subsection{5.6 Die Z4-Erweiterung: erlaubte Dynamik und echte neue
Parameter}\label{die-z4-erweiterung-erlaubte-dynamik-und-echte-neue-parameter}

Der neueste Nachtrag zeigt außerdem, dass die Erhaltung von \(N\)
stärker ist als die innere Symmetrie allein verlangt. Ein auf jedem Grad
konstanter ganzzahliger Lift der betrachteten Z4-Klammern müsste

\[2q_1=q_2,\qquad q_1+q_3=0,\qquad2q_3=q_2\]

erfüllen. Die Koeffizientenmatrix besitzt Determinante vier; ganzzahlig
bleibt nur der triviale Lift. Das widerlegt keine anderen Cartanladungen
und setzt die E8-Lieklammer nicht mit primitiven CAR-Operatoren gleich.
{[}S074, B4{]}

Die reelle symmetrische invariante Bosonpaarung \(\eta\) mit
\(\eta^2=I\) erlaubt am selben Tensor zusätzlich

\[B_+=\frac12b^\dagger\eta b^\dagger,\qquad
R_+=b^\dagger\eta P^\dagger.\]

Beide Operatoren sind \(G\)-invariant, ändern \(N\) um vier und erhalten
\(N\bmod4\). Innerhalb der ausdrücklich begrenzten Klasse
normalgeordneter, fermionparitätsgerader, \(G\)-invarianter Polynome vom
Grad höchstens drei ergibt sich

\[\begin{aligned}
H_{\mathrm{ext}}={}&c+\varepsilon N_f+\Delta N_b
+\kappa B_++\overline\kappa B_-\\
&+g\sum_A b_A^\dagger P_A
+\overline g\sum_A P_A^\dagger b_A
+\lambda R_++\overline\lambda R_-.
\end{aligned}\]

Die Quellenprüfung der vollen Generatoren zeigt die Zulässigkeit. Sie
leitet die Kopplungen nicht aus dem Ursprung ab. Die nativen
Grundzustands- und Polschranken gelten für diese Erweiterung nicht
automatisch. \(N\bmod4\) ist außerdem bereits durch das Zentrum von
\(SU(4)\) realisiert. Der Übergang von neun auf acht kontinuierliche
diagonale Phasenerhaltungen bedeutet nicht, dass die volle innere Gruppe
verschwände. Ein zusätzlicher \(\mu N\)-Term bleibt \(G\)- und
Z4-invariant; die Auswahlfrage bleibt bestehen. Für \(\Delta>|\kappa|\)
lässt sich die Energie nach unten beschränken, doch daraus folgt noch
kein eindeutiger neuer Grundzustand. {[}S074, B4.1{]}

\textbf{Wann ist die Erweiterung nur eine andere Beschreibung?} In der
reellen, phasengleichen Familie führt die kanonische Transformation

\[b=c\cosh r+\eta c^\dagger\sinh r\]

zu

\[\begin{aligned}
\Delta'&=\Delta\cosh2r+\kappa\sinh2r,&
\kappa'&=\Delta\sinh2r+\kappa\cosh2r,\\
g'&=g\cosh r+\lambda\sinh r,&
\lambda'&=g\sinh r+\lambda\cosh r.
\end{aligned}\]

Für \(g^2>\lambda^2\) lassen sich die beiden zusätzlichen Terme durch
diese uniforme Transformation gleichzeitig entfernen genau dann, wenn

\[\boxed{\mathcal I=\kappa(g^2+\lambda^2)-2\Delta g\lambda=0.}\]

Die additive Konstante ist \(30(\Delta'-\Delta)\). Bei
\(\mathcal I\ne0\) ist diese spezielle Reduktion ausgeschlossen; ein
Satz über alle denkbaren verborgenen Symmetrien folgt nicht. Die
einfache Familie \(\lambda=g\), \(\kappa=0\) entgeht einer solchen
endlichen Reduktion und ist ein konkreter Kandidat für weitere
Forschung. Bei \(\varepsilon=0\) besitzen beide kubischen Kanäle
dieselbe freie Frequenz \(\Delta\); ein bloßes Argument mit schneller
und langsamer Rotation erklärt ihr unterschiedliches Behalten nicht.
{[}S074, B4.2{]}

Dieser Nachtrag verbessert die fundamentale Fragestellung erheblich:
Zuerst müssen physikalisch verschiedene Dynamiken von bloßen
Variablenwechseln getrennt werden. Erst danach kann eine Ursprungsregel
sinnvoll auswählen. Die Auswahl erfolgt also auf Klassen gleicher
vollständiger Prozesse, was unmittelbar an Abschnitt 8 anschließt.

\section{6. Der neueste Stand v1.6.7: die richtige kleine
Rechenkette}\label{der-neueste-stand-v1.6.7-die-richtige-kleine-rechenkette}

\subsection{6.1 Vier Singuletts auf der zweiten
Bosonstufe}\label{vier-singuletts-auf-der-zweiten-bosonstufe}

Sei \(F\) der vollständig gefüllte Fermionzustand ohne Bosonen. Bei
\(k\) Bosonen gibt es im \(N=64\)-Sektor \(2k\) Löcher. Der invariante
Teil liegt in

\[\mathcal S_k=
[\Lambda^{2k}(\overline{16\otimes4})\otimes
\operatorname{Sym}^k(10\otimes6)]^G.\]

Die ersten Dimensionen sind \(1,1,4\). Die vierte Richtung bedeutet
keinen vierten Grundzustand, sondern einen größeren Raum möglicher
Singulettkombinationen.

Auf der Bosonseite gilt

\[\operatorname{Sym}^2(10\otimes6)=
(1,1)\oplus(54,1)\oplus(1,20')\oplus(54,20')\oplus(45,15).\]

Der skalare \((1,1)\)-Typ hat keinen passenden Beitrag in der
betreffenden Fermionstufe. Die vier übrigen gemeinsamen Typen liefern
die vollständige Singulettbasis. Für \(v_2=T_+^2F\) lauten die
Normquadrate und Kopplungen:

{\def\LTcaptype{none} % do not increment counter
\begin{longtable}[]{@{}lrr@{}}
\toprule\noalign{}
Typ & \(\|P_Rv_2\|^2\) & quadrierte Kopplung/\(g^2\) \\
\midrule\noalign{}
\endhead
\bottomrule\noalign{}
\endlastfoot
\((54,1)\) & 17280 & 36 \\
\((1,20')\) & 7680 & 16 \\
\((54,20')\) & 241920 & 504 \\
\((45,15)\) & 172800 & 360 \\
Summe & 439680 & 916 \\
\end{longtable}
}

Die vollständige Kompression bis einschließlich dieser Stufe ist

\[
P_{\le2}HP_{\le2}\simeq
\begin{pmatrix}
0&g\sqrt{480}&0\\
g\sqrt{480}&\Delta&g\sqrt{916}\\
0&g\sqrt{916}&2\Delta
\end{pmatrix}\oplus2\Delta I_3.
\]

Sie ist exakt bestimmt, aber nicht invariant: Übergänge zur nächsten
Bosonstufe bleiben. Ihre Eigenwerte sind Variationsdaten, nicht
automatisch Eigenwerte des vollständigen Operators. {[}S068, §3{]}

